\documentclass[10pt,a4paper]{amsart}
\usepackage{amsmath,amssymb,amsthm,graphicx}
\usepackage{verbatim}
\newenvironment{reference}{\comment}{\endcomment}
%\newenvironment{reference}{}{}
\newenvironment{slogan}{\comment}{\endcomment}
\newenvironment{history}{\comment}{\endcomment}

% For commutative diagrams we use Xy-pic
\usepackage[all]{xy}

% We use 2cell for 2-commutative diagrams.
\xyoption{2cell}
% Theorem environments.
%
\theoremstyle{plain}
\newtheorem{theorem}[subsection]{Theorem}
\newtheorem{proposition}[subsection]{Proposition}
\newtheorem{lemma}[subsection]{Lemma}

\theoremstyle{definition}
\newtheorem{definition}[subsection]{Definition}
\newtheorem{example}[subsection]{Example}
\newtheorem{exercise}[subsection]{Exercise}
\newtheorem{situation}[subsection]{Situation}

\theoremstyle{remark}
\newtheorem{remark}[subsection]{Remark}
\newtheorem{remarks}[subsection]{Remarks}

\numberwithin{equation}{subsection}

% Macros
%
\def\lim{\mathop{\mathrm{lim}}\nolimits}
\def\colim{\mathop{\mathrm{colim}}\nolimits}
\def\Spec{\mathop{\mathrm{Spec}}}
\def\Hom{\mathop{\mathrm{Hom}}\nolimits}
\def\Ext{\mathop{\mathrm{Ext}}\nolimits}
\def\SheafHom{\mathop{\mathcal{H}\!\mathit{om}}\nolimits}
\def\SheafExt{\mathop{\mathcal{E}\!\mathit{xt}}\nolimits}
\def\Sch{\mathit{Sch}}
\def\Mor{\mathop{\mathrm{Mor}}\nolimits}
\def\Ob{\mathop{\mathrm{Ob}}\nolimits}
\def\Sh{\mathop{\mathit{Sh}}\nolimits}
\def\NL{\mathop{N\!L}\nolimits}
\def\CH{\mathop{\mathrm{CH}}\nolimits}
\def\proetale{{pro\text{-}\acute{e}tale}}
\def\etale{{\acute{e}tale}}
\def\QCoh{\mathit{QCoh}}
\def\Ker{\mathop{\mathrm{Ker}}}
\def\Im{\mathop{\mathrm{Im}}}
\def\Coker{\mathop{\mathrm{Coker}}}
\def\Coim{\mathop{\mathrm{Coim}}}

% Boxtimes
%
\DeclareMathSymbol{\boxtimes}{\mathbin}{AMSa}{"02}

%
% Macros for moduli stacks/spaces
%
\def\QCohstack{\mathcal{QC}\!\mathit{oh}}
\def\Cohstack{\mathcal{C}\!\mathit{oh}}
\def\Spacesstack{\mathcal{S}\!\mathit{paces}}
\def\Quotfunctor{\mathrm{Quot}}
\def\Hilbfunctor{\mathrm{Hilb}}
\def\Curvesstack{\mathcal{C}\!\mathit{urves}}
\def\Polarizedstack{\mathcal{P}\!\mathit{olarized}}
\def\Complexesstack{\mathcal{C}\!\mathit{omplexes}}
% \Pic is the operator that assigns to X its picard group, usage \Pic(X)
% \Picardstack_{X/B} denotes the Picard stack of X over B
% \Picardfunctor_{X/B} denotes the Picard functor of X over B
\def\Pic{\mathop{\mathrm{Pic}}\nolimits}
\def\Picardstack{\mathcal{P}\!\mathit{ic}}
\def\Picardfunctor{\mathrm{Pic}}
\def\Deformationcategory{\mathcal{D}\!\mathit{ef}}
\usepackage{cleveref}
\crefname{originalreference}{original reference}{original references}
\Crefname{originalreference}{Original reference}{Original references}
\makeatletter
\newlabel{dpa-lemma-kernel}{{07H2}{source}{}{sourcepointer-1}{}}
\newlabel{dpa-lemma-kernel@cref}{{[originalreference][0][]07H2}{[1][0][]source}{}{}{}}
\makeatother
\title{Divided power envelope of a quotient ring}
\author{The Stacks Project Authors}
\date{stacks-project:crystalline.tex:local-snapshot-20261012; GNU FDL 1.2 or later}
\begin{document}
\maketitle
\noindent\textit{Editorial note.} All delivered mathematical text is the author's own, copied verbatim from the pinned Stacks Project chapter file crystalline.tex; no statement or proof step was rewritten, reordered or removed.\par\medskip
\noindent\textit{Editorial note.} Span recomputation by label pairing (LANE\_BRIEF 10.2): in this chapter the author's proof environment follows the statement environment as a standalone block; the spec covers the complete environment pair (statement span, proof span) as recorded in the pool candidate.\par\medskip
\ifdefined\hypertarget\hypertarget{sourcepointer-1}{}\fi
\noindent\textit{Source reference.} Original reference 07H2: sources/stacks-chow/tags.txt, source lines 9537--9537. Cross-chapter standard prerequisite cited by the author inside the counted span; it is used as a pointer and its own proof is not part of this item.\par\medskip
\begin{lemma}
\label{lemma-divided-power-envelop-quotient}
Let $(A, I, \gamma)$ be a divided power ring.
Let $\varphi : B' \to B$ be a surjection of $A$-algebras with kernel $K$.
Let $IB \subset J \subset B$ be an ideal. Let $J' \subset B'$
be the inverse image of $J$. Write
$D_{B', \gamma}(J') = (D', \bar J', \bar\gamma)$.
Then $D_{B, \gamma}(J) = (D'/K', \bar J'/K', \bar\gamma)$
where $K'$ is the ideal generated by the elements $\bar\gamma_n(k)$
for $n \geq 1$ and $k \in K$.
\end{lemma}

\begin{proof}
Write $D_{B, \gamma}(J) = (D, \bar J, \bar \gamma)$.
The universal property of $D'$ gives us a homomorphism $D' \to D$
of divided power algebras. As $B' \to B$ and $J' \to J$ are surjective, we
see that $D' \to D$ is surjective (see remarks above). It is clear that
$\bar\gamma_n(k)$ is in the kernel for $n \geq 1$ and $k \in K$, i.e.,
we obtain a homomorphism $D'/K' \to D$. Conversely, there exists a divided
power structure on $\bar J'/K' \subset D'/K'$, see
Divided Power Algebra, Lemma \ref{dpa-lemma-kernel}.
Hence the universal property of $D$ gives an inverse $D \to D'/K'$ and we win.
\end{proof}

\end{document}
